\chapter{GEO --- Generative Engine Optimisation}
\label{ch:geo}

\section{Applicability and scope}

Generative Engine Optimisation concerns whether a product is represented and
cited accurately by generative systems: large language models and the assistants
built on them, including code-generating assistants. For a developer tool, GEO
matters because a common way a developer learns about a library is to ask an
assistant ``how do I scaffold an Expo login screen?''.

There is no project-owned website to optimise, so classic GEO tactics
(rich pages, citations, entity markup) are \textbf{NOT APPLICABLE}. What is
applicable is the \emph{corpus footprint}: whether the project's own material is
present, self-consistent and machine-consumable where generative systems look.

\section{Corpus footprint}

\begin{longtable}{@{}L{4.4cm}L{2.8cm}L{7cm}@{}}
\caption{Sources a generative system could learn from}\\
\toprule
\textbf{Source} & \textbf{Availability} & \textbf{Content} \\
\midrule
\endfirsthead
\toprule
\textbf{Source} & \textbf{Availability} & \textbf{Content} \\
\midrule
\endhead
Public GitHub repository & Public (master branch) & Source, recipes, templates,
workflows, READMEs. \\
npm packages & Public & \texttt{@expo-snips/core}, \texttt{@expo-snips/cli}
metadata and (if a README is published) documentation. \\
VS Code Marketplace listing & Public (when listed) & Extension metadata, README,
keywords. \\
Recipe registry JSON & Public & Structured catalog with descriptions and tags. \\
Answering-assistant context & Not controlled & Unknown whether any model has been
trained on or cites this project. \\
\bottomrule
\end{longtable}

\section{Generative-readiness properties}

\begin{itemize}
  \item \textbf{Structured, self-describing data:} the recipe catalog is JSON
        Schema-conformant, with \texttt{id}, \texttt{name}, \texttt{description}
        and \texttt{tags} per recipe. This is precisely the shape a generative
        system can reason over, and it is a genuine strength.
  \item \textbf{Machine-readable API documentation:} the \texttt{@expo-snips/core}
        package ships TypeScript declarations, and the generated
        \texttt{registry-data.ts} is typed, giving models accurate signatures.
  \item \textbf{Versioned, dated history:} the extension CHANGELOG is dated and
        factual.
  \item \textbf{No telemetry or private corpus:} nothing about the tool's usage is
        captured, so there is no private signal to leak or misrepresent.
\end{itemize}

\textbf{Status:} IMPLEMENTED.

\section{Representation risks}

\begin{longtable}{@{}L{1.5cm}L{4.6cm}L{3.4cm}L{4.7cm}@{}}
\caption{GEO findings}\\
\toprule
\textbf{ID} & \textbf{Finding} & \textbf{Evidence} & \textbf{Recommendation} \\
\midrule
\endfirsthead
\toprule
\textbf{ID} & \textbf{Finding} & \textbf{Evidence} & \textbf{Recommendation} \\
\midrule
\endhead
GEO-01 & Inconsistent identity across sources: the extension publishes under
publisher \texttt{ukown812} and repository \texttt{github.com/unkown812/expo-snips},
while some README links point at \texttt{github.com/expo-snips/expo-snips}. A
generative system may conflate these into a wrong citation. & Extension
\texttt{package.json} vs README links. & Canonicalise one repository URL and apply
it everywhere. \\
GEO-02 & Name collision risk: the extension is published as \texttt{expo-snippets},
the CLI as \texttt{@expo-snips/cli}, and the product as ``Expo Snips''. Models may
treat these as three products. & \texttt{package.json} names. & Keep the
``Expo Snips'' brand string adjacent to each package name in its description. \\
GEO-03 & Stale prose in \texttt{ProjectDescription.md} can propagate an
incorrect fact (recipe count, variant support) into a generated answer. &
\texttt{documents/ProjectDescription.md}. & Correct the document; keep a single
source of truth for counts. \\
GEO-04 & No license clarity for the extension beyond ``Private License'', which a
generative system cannot interpret as a known SPDX license. & Extension
\texttt{package.json} \texttt{license} field. & Use an explicit, recognisable
license expression if the intent allows. \\
GEO-05 & No canonical, quotable ``what is this'' one-liner shared across sources. &
Three differing descriptions. & Publish one canonical sentence and reuse it. \\
\bottomrule
\end{longtable}

\section{Why GEO matters for this product specifically}

The product's value proposition --- ``generate a complete, correct screen for a
given stack'' --- is exactly the kind of question developers pose to coding
assistants. If a model has never seen the project or cites it inconsistently, the
assistant will instead suggest generic manual boilerplate, which is the
alternative the tool exists to replace. GEO is therefore economically meaningful
even without a website: the goal is that a generative system, asked for an Expo
screen, can name the tool and be correct about what it does.

\section{Summary}

\begin{table}[H]
\centering
\caption{GEO posture}
\begin{tabular}{@{}L{4.6cm}L{9.6cm}@{}}
\toprule
\textbf{Dimension} & \textbf{Assessment} \\
\midrule
Website/entity optimisation & NOT APPLICABLE (no site). \\
Structured corpus & IMPLEMENTED (schema-valid JSON, typed declarations). \\
Identity consistency & GAP (publisher/repo/brand mismatch). \\
Factual consistency & GAP (stale description document). \\
Model citations & UNKNOWN (not observable from the repository). \\
\bottomrule
\end{tabular}
\end{table}

\textbf{Status:} PARTIAL / RECOMMENDED (canonical identity, accurate descriptions).
